\part{Lô-gích \\ và phương pháp lập luận diễn dịch}

\renewcommand{\thefigure}{\thepart.\arabic{figure}}
\setcounter{figure}{0}

\vspace*{\fill}

\begin{figure}[H]
   \centering
   \includegraphics[width=0.8\textwidth]{images/nen_tang_lo_gich.jpg}
   \caption{Hình ảnh bàn cờ vua. \\ Nguồn: https://www.pexels.com/photo/brown-chess-pieces-on-brown-wooden-chess-board-5491882/}
   \label{fig:nen_tang_lo_gich}
\end{figure}

Giống như một ván cờ vua hoạt động dựa trên những quân cờ và luật di chuyển nghiêm ngặt, hệ thống tri thức của chúng ta cũng được vận hành trên những quy tắc suy luận rất chặt chẽ. Lô-gích học chính là ``luật chơi'' của tư duy; từ một vài tiên đề cơ bản ban đầu, chúng ta có thể kiến tạo nên vô vàn định lí phức tạp, giống như cách vô số thế cờ thiên biến vạn hóa được sinh ra chỉ từ một bàn cờ giới hạn.

\renewcommand{\thefigure}{\thechapter.\arabic{figure}}
\setcounter{figure}{0}

\vspace*{\fill}
\clearpage
\input{chapter/nen_tang_lo_gich/giai_tich_menh_de.tex}
\input{chapter/nen_tang_lo_gich/giai_tich_vi_tu_bac_nhat.tex}